% Combined, plain-language summary of Chapter 6 (Methodology) and Chapter 7
% (Project Structure \& Implementation). Written so a non-specialist reader
% can follow the whole system without either chapter's full technical
% detail. Every fact below is the same fact reported in Chapters 6 and 7 --
% nothing here is new, and nothing is repeated twice within this document.
% Where the two chapters described the same thing from two different
% angles, that description appears here exactly once, in the place it
% belongs (matching chapters/WRITING_RULES.md \S4's "no repetition" rule).
%
% Not part of the numbered chapter sequence -- a standalone companion
% document. \input{} it directly where wanted, or read on its own.

\section*{How the Pipeline Works --- a Plain-Language Guide}

\subsection*{1. What Problem This Part of the Thesis Solves}

The thesis treats ``is a tipping point coming, and what kind?'' as a
classification problem: feed a short stretch of a time series into a
model, and get back a guess at one of four labels -- Fold, Hopf,
Transcritical, or ``nothing happening'' (Null) -- instead of computing one
hand-picked statistic like variance and eyeballing whether it looks
worrying.

Two very different kinds of data go through this same system: a large
synthetic dataset with known, correct labels (used for training and for
checking how good a model really is), and the real Mediterranean
sediment-core measurements, which have no known label for bifurcation
type at all. The whole system is built so both kinds of data flow through
the \emph{same} pipeline, up to the point where the results are scored --
only the scoring step treats them differently, because only one of them
has a known answer to check against.

\subsection*{2. The Pipeline, Step by Step}

\begin{enumerate}
    \item \textbf{Raw data comes in.} Either a synthetic simulated series,
    or a real measurement of an element (molybdenum, Mo, or uranium, U)
    from a sediment core. Both are just a plain list of numbers at this
    point.
    \item \textbf{The data is cleaned and reshaped.} Each series is cut
    into fixed-length windows and each window is rescaled by dividing by
    its own average size, so that a window with big swings and one with
    small swings both end up on a comparable scale. Nothing is smoothed
    onto an artificial evenly-spaced timeline -- the code keeps the
    data's own real spacing, because smoothing it out would quietly wipe
    away the very signal (rising noisiness) the whole method is trying to
    detect, matching \citet{bury2021}'s own choice for the same reason.
    This step is handled by \texttt{src/dataset\_loader.py} for the
    synthetic data and \texttt{src/rolling\_window.py} for the real core
    data.
    \item \textbf{The path splits in two, depending on which model
    family is about to read the window} -- not depending on which of the
    two data sources it came from. A synthetic window and a real window
    go through exactly the same fork:
    \begin{itemize}
        \item The 22 ``classical'' models (statistical/kernel-based
        methods, not neural networks) get an enriched version of the
        window with five numbers per time point instead of one: the raw
        value, how noisy the recent stretch has been (rolling variance),
        how strongly each point resembles the one before it (rolling
        lag-1 autocorrelation), the lopsidedness of the recent values
        (rolling skewness), and how much the noise level has grown
        relative to where the window started (the ``variance growth
        ratio''). This step lives in \texttt{src/ews\_augmenter.py}.
        \item The 7 deep-learning models skip that step and just read
        the single raw, rescaled number per time point. They have to
        work out any useful patterns themselves, straight from the raw
        shape.
    \end{itemize}
    \item \textbf{A model makes a prediction.} Every one of the 29
    models -- all 7 deep-learning ones and all 22 classical ones --
    exposes the same \texttt{predict\_proba()} function, so the rest of
    the code never has to know or care which kind of model it is talking
    to. For a neural network, this is a standard softmax turning the
    network's raw output into four probabilities that add up to 1. For
    most of the classical models it is just that model's own built-in
    probability output. For one specific family of classical models (the
    ROCKET family, which only produce a raw decision score rather than a
    probability), the code applies the same softmax conversion by hand.
    Either way, the result leaving this step always looks the same: four
    probabilities, one per class.
    \item \textbf{The result is scored -- differently, depending on the
    data source.} For a synthetic window, the true label is known, so
    the model's four probabilities can be checked against ground truth
    directly. For a real core segment, the true bifurcation type is
    \emph{not} known, so only a narrower question can be answered: did
    the model detect that \emph{something} was approaching at all?
    Section~4 (Section~\ref{sec:evaluation_metrics} in the full
    Methodology chapter) explains exactly how each of those two checks
    works.
\end{enumerate}

Training and prediction never touch the real Mediterranean data at the
same time -- every model is trained entirely on the synthetic data first,
and Section~5 below covers exactly when and why its weights stop
changing after that.

\subsection*{3. Why Two Model Families, Not Just One}

The original method this thesis builds on used a single neural network
\citep{bury2021}. This thesis benchmarks both deep-learning models
\emph{and} 22 classical models, for a concrete, checkable reason: on
general-purpose time-series benchmarks, fast classical methods (kernel-
and interval-based ones especially) are often just as accurate as neural
networks, and considerably cheaper to train -- MiniRocket in particular
is reported as up to 75 times faster than the original ROCKET method
while keeping essentially the same accuracy \citep{dempster2021minirocket},
and it needs no GPU at all. Whether that general pattern also holds for
this specific bifurcation-classification task was never tested before
this thesis, so it was treated as an open question to actually answer,
not assumed.

\subsubsection*{The models themselves}

\textbf{Deep learning (7 models).} Three represent the three broad
design families in the roster and are explained in full architectural
detail in the Methodology chapter:

\begin{itemize}
    \item \textbf{LSTM} -- a plain recurrent network (not the
    architecture the original paper used; this one follows
    \citet{ma2025sdml}'s design instead), reading the raw series one
    step at a time through a memory cell that avoids the ``forgetting''
    problem simple recurrent networks have on long sequences.
    \item \textbf{CNN-LSTM} -- the original paper's own architecture
    \citep{bury2021}: a convolutional layer first picks out local
    patterns (a short spike in noisiness, a brief pattern in how points
    relate to their neighbours) and compresses them, and only that
    compressed summary is handed to an LSTM, so the LSTM never has to
    track every raw noisy point itself.
    \item \textbf{InceptionTime} \citep{fawaz2020inceptiontime} -- runs
    several convolutional filters of different lengths in parallel over
    the same window, so the network reads the series at several time
    scales at once (useful because a Hopf bifurcation's oscillation is
    fast while a Fold's slow build-up is gradual).
\end{itemize}

Four more were added to the roster and each closely reproduces a
published architecture, checked layer-by-layer against a reference
implementation: \textbf{ResNet} \citep{wang2017resnet} (residual blocks
with a shortcut connection around each one), \textbf{TCN}
\citep{bai2018tcn} (a stack of convolutions that each see a
progressively wider stretch of the series, so the last one covers the
whole window), \textbf{RNN-FCN} \citep{karim2018lstmfcn} (an LSTM branch
and a convolutional branch combined, where the LSTM is deliberately
shown the whole series as one long step rather than many short ones),
and \textbf{PatchTST} \citep{nie2023patchtst} (splits the series into
short patches and treats each patch as one token for a transformer, the
same general kind of architecture behind modern language models).

\textbf{Classical, non-neural-network models (22 models),} surveyed
comprehensively by \citet{bagnall2017bakeoff} and grouped into three
families:

\begin{itemize}
    \item \textbf{Kernel-based (the ROCKET family)}
    \citep{dempster2020,dempster2021minirocket,tan2022multirocket,
    middlehurst2021hivecote2} -- convolve the series with a large bank
    of fixed or randomly-generated filters and, for each filter, just
    count what fraction of the output is positive; feed that into a
    simple linear classifier. Fast, and often as accurate as much
    heavier methods.
    \item \textbf{Dictionary-based} \citep{schafer2023weasel2,
    senin2013saxvsm} -- turn each window into a short discrete ``word''
    describing its shape, then classify based on which words show up how
    often, similar in spirit to how spam filters count words in an
    email.
    \item \textbf{Interval- and feature-based} \citep{deng2013tsf} --
    compute simple statistics (mean, slope, standard deviation, and
    similar) over many different sub-stretches of the window and combine
    them with a standard decision forest.
\end{itemize}

\subsection*{4. How a Result Is Measured}

A model could just be judged by raw accuracy (how often its single top
guess is right), but that number hides how the model behaves at other
confidence thresholds, and for an early-warning system the acceptable
trade-off between missing a real transition and raising a false alarm is
itself a judgement call, not a fixed 50\% cut-off. This thesis instead
uses AUC -- the area under the ROC curve, a standard measure that
summarises performance across every possible threshold at once, from 0.5
(no better than a coin flip) to 1.0 (perfect). It is used here
specifically because comparing 29 very different models fairly needs a
measure that does not depend on every model happening to be equally
well-calibrated at the exact same cut-off, and AUC does not depend on
that.

Two different versions of this same AUC calculation are used, matching
the two different questions from Section~2 above: on the synthetic data,
where the correct label is known, each of the four classes is scored
against the other three and the four scores are averaged -- this
measures how well a model tells each specific bifurcation type apart
from everything else. On the real core data, where the type is not
known, only the narrower ``was a transition detected at all'' question
is scored, against simulated no-transition control data.

\subsection*{5. How Training Works, and When the Model Stops Changing}

The synthetic dataset -- the only one with known correct labels -- is
split once into three fixed parts: 95\% for training, 4\% for adjusting
the model along the way (deciding when to stop training and when to slow
the learning rate down), and 1\% held back untouched to give one honest
accuracy number per model at the end. This is a single clean split, not
repeated in different ways.

Every deep-learning model is trained with the same loss function
(categorical cross-entropy, the standard way to penalise a wrong
probability guess across several classes) and the Adam optimiser
\citep{kingma2015adam}. Most of them also use a scheduler that
automatically halves the learning rate once progress stalls; one model
(the plain LSTM) is the exception and trains at one fixed learning rate
throughout. Batch size, exactly how patient training is before stopping,
and similar per-model settings are each tuned individually rather than
fixed to one value for the whole roster.

Once training and tuning are done, every model's weights are frozen and
saved to disk -- deep-learning weights in one standard PyTorch format
\citep{paszke2019pytorch}, classical-model objects in another standard
Python format -- and never updated again. This is the point where the
real Mediterranean cores enter the picture for the first time: a frozen
model is simply run on them, with no further learning of any kind, so
the real-data result is a genuine test of what transferred from
synthetic training, not a retrain on the target domain. Section~6 covers
exactly how that real-data run works.

\subsection*{6. How the Real Mediterranean Data Is Actually Evaluated}

A trained, frozen model is run against a real sediment core through
\texttt{testing/evaluate.py}. Rather than producing one single verdict
for an entire core, the code takes 40 evenly-spaced windows stretching
from the start of the usable pre-transition record up to the transition
itself, and predicts once per window -- this gives a whole trajectory of
predictions approaching the transition, not a single answer. Each
window's four-class probabilities, the derived ``how likely is a
transition'' number, and the classical rolling variance and
autocorrelation for that same window are all written out together, and
every summary table and chart in the results chapter is ultimately built
from these saved per-window numbers.

\subsection*{7. How the Code Itself Is Organised}

The codebase is deliberately organised around a small number of files
that each do one job, connected by plain function calls -- not around a
hierarchy of custom classes. Loading data, building the extra input
channels, defining each model, training, and evaluating are each their
own file with one or two top-level functions; there is no single
``Trainer'' or ``Model'' class managing everything, and no separate
script just for running the real-data evaluation -- it is the same
evaluation file used for the synthetic test data, with a different flag.
Adding a brand-new model to the roster means adding one new file and one
new entry to a settings file, not extending an inheritance chain. This
was a deliberate choice, checked against \citet{bury2021}'s own published
code \citep{burygithub}, which is organised the same simple,
function-based way.

None of the 22 classical models were coded from scratch -- each one is
built directly from an existing, well-tested open-source library (mainly
\texttt{aeon} \citep{aeon2024}, with \texttt{pyts} \citep{pyts2020}
covering the handful of models \texttt{aeon} does not provide), so this
thesis's own code is responsible for wiring the right library call
together with the right settings, not for reimplementing the underlying
algorithm. For the ROCKET-family models specifically, the underlying
library pairs its kernel features with a ridge-regression classifier
\citep{pedregosa2011sklearn}, not code written for this thesis.

\subsection*{8. Where the Experiments Actually Ran}

Training and testing 29 models at two different window lengths adds up
to far more computation than a single computer could reasonably do. The
experiments ran on a university computing cluster, using a standard job
scheduler, SLURM \citep{yoo2003slurm}, to launch one job per model
automatically rather than by hand. Deep-learning jobs were given a
graphics-card partition, since their bottleneck is the heavy matrix
arithmetic behind training a neural network; the classical models --
which involve no such training step, just fast feature extraction --
were instead given many ordinary CPU cores and no graphics card, matching
what each family of methods actually needs to run efficiently.

\bigskip
\noindent\emph{Every number, file name, and function name above matches
exactly what Chapter~\ref{ch:methodology} and Chapter~\ref{ch:implementation}
report; see those chapters (and their matching
\texttt{chapter\_6\_verification.md} / \texttt{chapter\_7\_verification.md}
logs) for the full technical detail, exact architectural specifications,
and source-level verification behind each claim.}
